\chapter{先说重要的}

本文在 TeX Live 软件环境下, 使用 XeLaTeX 引擎, 可以顺利编译.
TeX Live 的安装方法, 请移步查看: \url{http://aff.whu.edu.cn/huangzh/}.


\section{文档具体使用步骤}

\begin{description}
  \item[Step 1]  进入 data 文件夹,  打开 whusetup.tex, innovations.tex,  cnabstract.tex,  enabstract.tex, nomenclature.tex, achievements.tex, thanks.tex 这几个文档,分别填写题目与作者等等信息、创新点、中文摘要、英文摘要、缩略词、成果列表以及致谢.

  \item[Step 2]  在data文件夹中新建chap文件，书写正文，并打开主文档 PhdThesis.tex, 添加chap文件.

  \item[Step 3]  使用 XeLaTeX{} 编译. 具体见 \cref{sec-compile} 节.

\end{description}


\section{编译的方法}\label{sec-compile}

默认使用 XeLaTeX 编译, 直接生成~pdf 文件.

若另存为新文档, 请确保文档保存类型为 \verb|:UTF-8|. 当然目前很多编辑器默认文字编码为 UTF-8.
WinEdt 9.0 之后的版本都是默认保存为 UTF-8 的.


\section{打印的问题}
\begin{enumerate}[label=(\arabic*)]
  \item  论文双面打印.
  \item  打印时留意不要缩小页面或居中. 即页面放缩方式应该是``无''(Adobe Reader XI 是选择``实际大小'').
           有可能页面放缩方式默认为``适合可打印区域'', 会导致打印为原页面大小的 $97\%$.
           文字不要居中打印, 是因为考虑到装订, 左侧的空白留得稍多一点(模板已作预留).
  \item  遗留问题: 封面需要打印部重新制作. 比如博士论文要打印书脊, 校内打印部通常有现成的模板.
           我们自己做的封面, 打印部不一定好用.
\end{enumerate}
%如果不是彩色打印机, 请在打印时, 选择将彩色打印为黑白, 否则彩色文字打出的墨迹会偏淡.




\section{提交电子文档}

向武汉大学图书馆提交电子文档, 需单独编译文件: 在文档选项中添加 forlib.

原因: 图书馆要求提交的电子文档不能有空白页、彩色文字.

({\kaishu 这个要求使我在编制模板时遇到了一点问题: 这会导致电子版与纸质版的页码不一致. 博士论文每章的起始页默认在奇数页, 这会不可避免地出现空白页.})


\vfill

本文档下载更新地址: \url{http://aff.whu.edu.cn/huangzh/}. 使用之前, 请移步查看是否有更新.

问题反馈及建议, 请联系: huangzh@whu.edu.cn.
